\chapter{Discussion}
\label{ch:discussion}

This chapter steps back from the individual results to their broader meaning: what the two parts jointly
say about intent-driven network control (\cref{sec:disc:interpretation}), why the findings matter beyond
the specific systems (\cref{sec:disc:implications}), and what the results cannot tell us
(\cref{sec:disc:limitations}).

\section{Interpretation}
\label{sec:disc:interpretation}

Read together, the two parts tell a consistent story about \emph{when} additional control machinery earns
its keep.

\paragraph{Intent conditioning turns a family of policies into one.} Part~I's central lesson is that the
right way to serve many objectives is not many models but one model that takes the objective as input --
provided the objective is allowed to \emph{modulate} the decision rather than merely accompany it. The
concatenation-versus-FiLM contrast (\cref{sec:p1design:film}) is the crisp version of this point: a
conditioning signal that enters symmetrically across the alternatives cancels under the comparison that
picks the action, so the policy can look conditioned while behaving identically. Feature-wise modulation
avoids this because it changes the \emph{relative} salience of path features, and thus which path wins. This
is a design principle, not a SCION-specific trick: any comparison-based selector conditioned on an
objective should modulate its per-alternative features.

\paragraph{The value of multipath control is a function of volatility.} Part~II's central lesson is that
learned multipath scheduling ranges from nearly worthless to decisive depending on the network. In the
static scenario, a good bitrate controller over a trivial split captured essentially all achievable QoE, and
the learned transport agent added little (\cref{sec:p2eval:results}). In the dynamic scenario, the same
app-only configuration collapsed to negative QoE while the learned split sustained a usable call. The
practical implication is that one should not evaluate multipath schedulers only in stable conditions -- that
is precisely the regime in which they look unnecessary -- and that the case for learned scheduling rests on
volatility, churn, and bursts, where static policies cannot adapt.

\paragraph{One principle, two layers.} Both parts instantiate the Mortise principle~\cite{Shen2026} --
intent as a first-class input that reshapes a network mechanism -- at successively deeper points of control.
Part~I reshapes \emph{which path} is chosen; Part~II reshapes both \emph{how the paths are used} and
\emph{what the application produces}. The progression suggests a general recipe: expose the application's
intent, and let it steer every controllable decision between the application and the wire.

\section{Implications}
\label{sec:disc:implications}

The results matter for how we build controllers on path-aware architectures. Prior learned path
selectors~\cite{Bourak2025} and learned congestion or scheduling controllers~\cite{Jay2019, Zhang2019,
Wu2020} each optimize a fixed objective; our work argues that on a path-aware substrate the objective should
be a knob, and that a single conditioned policy can replace a zoo of per-objective models. For real-time
media specifically, the joint bitrate-and-split result reinforces a theme from learned
ABR~\cite{Mao2017} -- that application and network decisions are better made together -- and extends it to
the multipath setting that SCION makes routine. More broadly, the two systems are evidence that SCION's
endpoint path choice is not merely a routing convenience but an enabler of application-aware optimization
that is awkward or impossible in a destination-based Internet.

\section{Limitations}
\label{sec:disc:limitations}

The results are bounded in ways that the analysis (\cref{ch:analysis}) details and that fairness requires
restating here. Both systems are evaluated in simulation, and Part~II's headline dynamic-scenario numbers
are from the mock backend pending NS-3 confirmation (\cref{sec:p2impl:gaps}). Intent in this thesis is
\emph{declared}, not \emph{inferred}: we assume the application can state its objective, and we do not study
how to recover a latent intent from traffic -- a natural but separate problem. The transport model is
TCP-like rather than true Multipath QUIC, and the VMAF surrogate is bitrate-dominated, so absolute QoE
figures should be read as relative comparisons rather than calibrated user scores. Finally, we have not
addressed online, safe learning against live traffic, which any deployment would require. None of these
undercut the two qualitative findings -- modulation-based conditioning works where concatenation fails, and
multipath control pays off in proportion to volatility -- but they do delimit how far the quantitative
results should be pushed.
